\par\Needspace{6\baselineskip}\medskip
\hypertarget{native-146-article106}{}\pdfbookmark[2]{106.}{article-106}
\begin{center}106.\end{center}
Bei der Behandlung der bemerkenswertesten Beziehungen für die Bewegung der Himmelskörper auf Kegelschnitten müssen wir den eleganten Ausdruck erwähnen, der die Zeit durch die große Halbachse, die Summe $r+r'$ und die Sehne zwischen den beiden Orten angibt. Für die Parabel scheint diese Formel zuerst \SourceLeaf{147}von dem berühmten \textsc{Euler} gefunden worden zu sein (Miscell. Berolin. T. VII p. 20). Er ließ sie jedoch später unbeachtet und dehnte sie auch nicht auf die Ellipse und die Hyperbel aus. Wer die Formel dem namhaften \textsc{Lambert} zuschreibt, irrt sich also, obwohl diesem Geometer das Verdienst zukommt, den in Vergessenheit geratenen Ausdruck selbstständig wiedergefunden und auf die übrigen Kegelschnitte ausgedehnt zu haben. Mehrere Geometer haben diesen Gegenstand bereits behandelt, doch aufmerksame Leser werden die folgende Darstellung für angebracht halten. Wir beginnen mit der elliptischen Bewegung.
\par
Zunächst bemerken wir, dass der um die Sonne beschriebene Winkel $2f$ (Art. 88, dem wir auch die übrigen Bezeichnungen entnehmen) kleiner als $360^0$ vorausgesetzt werden kann. Vergrößert man diesen Winkel um 360 Grad, so verlängert sich die Zeit offenbar um die Dauer eines Umlaufs, also um $\dfrac{a^{\frac32}\mathpunct{.}360^0}{k}=a^{\frac32}\times365{,}25$ Tage. Bezeichnen wir nun die Sehne mit $\rho$, so gilt offenbar
\[
\rho\rho=(r'\cos v'-r\cos v)^2+(r'\sin v'-r\sin v)^2
\]
\noindent und daher nach den Gleichungen VIII und IX in Art. 8
\[
\begin{aligned}
\rho\rho&=aa(\cos E'-\cos E)^2+aa\cos\varphi^2(\sin E'-\sin E)^2\\[.6em]
&=4aa\sin g^2(\sin G^2+\cos\varphi^2\cos G^2)\\[.6em]
&=4aa\sin g^2(1-ee\cos G^2).
\end{aligned}
\]
\noindent Wir führen einen Hilfswinkel $h$ ein, für den $\cos h=e\cos G$ gilt. Um jede Mehrdeutigkeit zu beseitigen, wählen wir $h$ zwischen 0 und $180^0$, sodass $\sin h$ positiv ist. Da $g$ ebenfalls zwischen denselben Grenzen liegt (denn würde $2g$ den Wert $360^0$ erreichen oder überschreiten, so würde die Bewegung um die Sonne einen ganzen Umlauf erreichen oder überschreiten), folgt aus der vorhergehenden Gleichung unmittelbar
\[
\rho=2a\sin g\sin h,
\]
\noindent sofern die Sehne als positive Größe aufgefasst wird. Da ferner
\[
r+r'=2a(1-e\cos g\cos G)=2a(1-\cos g\cos h),
\]
\noindent gilt, folgt offenbar mit $h-g=\delta$, $h+g=\epsilon$
\[
\begin{aligned}
\mathrm{1)}\qquad&r+r'-\rho=2a(1-\cos\delta)=4a\sin\tfrac12\delta^2\\[.6em]
\mathrm{2)}\qquad&r+r'+\rho=2a(1-\cos\epsilon)=4a\sin\tfrac12\epsilon^2.
\end{aligned}
\]
\par
Schließlich gilt
\[
kt=a^{\frac32}(2g-2e\sin g\cos G)=a^{\frac32}(2g-2\sin g\cos h),
\]
\noindent\SourceLeaf{148}oder
\[
\mathrm{3)}\qquad kt=a^{\frac32}(\epsilon-\sin\epsilon-(\delta-\sin\delta)).
\]
\par
Nach den Gleichungen 1 und 2 lassen sich also die Winkel $\delta$ und $\epsilon$ aus $r+r'$, $\rho$ und $a$ bestimmen. Damit lässt sich nach Gleichung 3 aus denselben Größen die Zeit $t$ bestimmen. Die Formel lässt sich nach Belieben auch in der folgenden Form schreiben.
\[
\begin{aligned}
kt=a^{\frac32}\Biggl\{&\operatorname{arc}\cos\frac{2a-(r+r')-\rho}{2a}\\[.6em]
&-\sin\operatorname{arc}\cos\frac{2a-(r+r')-\rho}{2a}\\[.6em]
&-\operatorname{arc}\cos\frac{2a-(r+r')+\rho}{2a}\\[.6em]
&+\sin\operatorname{arc}\cos\frac{2a-(r+r')+\rho}{2a}\Biggr\}\mathpunct{\cdot}
\end{aligned}
\]
\noindent Bei der Bestimmung der Winkel $\delta$, $\epsilon$ aus ihren Kosinuswerten bleibt eine Mehrdeutigkeit, die genauer zu betrachten ist. Offenbar muss $\delta$ zwischen $-180^0$ und $+180^0$ und $\epsilon$ zwischen 0 und $360^0$ liegen. Auch innerhalb dieser Grenzen scheint jeder der beiden Winkel zwei Bestimmungen und die sich ergebende Zeit damit vier Bestimmungen zuzulassen. Aus Gleichung 5 in Art. 88 haben wir jedoch
\[
\cos f\mathpunct{.}\BareRoot{rr'}=a(\cos g-\cos h)=2a\sin\tfrac12\delta\sin\tfrac12\epsilon:
\]
\noindent Nun ist $\sin\tfrac12\epsilon$ notwendigerweise positiv. Daraus folgt, dass $\cos f$ und $\sin\tfrac12\delta$ dasselbe Vorzeichen haben müssen. Daher ist $\delta$ zwischen 0 und $180^0$ oder zwischen $-180^0$ und 0 zu wählen, je nachdem, ob $\cos f$ positiv oder negativ ist, das heißt, ob die heliozentrische Bewegung $2f$ unter oder über $180^0$ liegt. Offenbar muss für $2f=180^0$ der Wert $\delta=0$ sein. Damit ist $\delta$ vollständig bestimmt. Die Bestimmung des Winkels $\epsilon$ bleibt dagegen notwendigerweise mehrdeutig, sodass sich für die Zeit stets \textit{zwei} Werte ergeben. Welcher davon der wahre ist, lässt sich nur entscheiden, wenn dies anderweitig feststeht. Der Grund dafür ist leicht einzusehen. Durch zwei gegebene Punkte lassen sich nämlich \textit{zwei} verschiedene Ellipsen legen, die beide einen Brennpunkt im selben gegebenen Punkt und dieselbe große Halbachse haben\footnote{Beschreibt man um den ersten Ort einen Kreis mit dem Radius $2a-r$ und um den zweiten einen weiteren mit dem Radius $2a-r'$, so liegt der andere Brennpunkt der Ellipse offenbar in einem Schnittpunkt dieser Kreise. Da es im allgemeinen Fall stets zwei Schnittpunkte gibt, ergeben sich zwei verschiedene Ellipsen.}. Die Bewegung vom ersten zum zweiten Ort vollzieht sich auf diesen Ellipsen offenbar in ungleichen Zeiten.
